\documentclass[aspectratio=169]{beamer}
\usetheme{Madrid}
\usecolortheme{seahorse}

\usepackage[spanish]{babel}
\usepackage[utf8]{inputenc}
\usepackage[T1]{fontenc}
\usepackage{amsmath, amssymb}
\usepackage{booktabs}
\usepackage{graphicx}
\usepackage{listings}
\usepackage{xcolor}

\title{Aplicacion de Metodos Numericos a un Modelo Empresarial}
\subtitle{Euler, Runge-Kutta de 2\textdegree y 4\textdegree orden}
\author{Nelson Guijarro}
\institute{Diplomado TICS}
\date{\today}

\definecolor{codebg}{RGB}{245,245,245}
\definecolor{codekw}{RGB}{0,90,156}
\definecolor{codecm}{RGB}{0,120,80}

\lstset{
  language=Python,
  backgroundcolor=\color{codebg},
  basicstyle=\ttfamily\footnotesize,
  keywordstyle=\color{codekw}\bfseries,
  commentstyle=\color{codecm},
  showstringspaces=false,
  frame=single,
  breaklines=true,
  tabsize=2
}

\begin{document}

\begin{frame}
  \titlepage
\end{frame}

\begin{frame}{Objetivo del trabajo}
\textbf{Proposito:} modelar y simular el crecimiento de ventas de una empresa con una EDO no lineal.

\vspace{0.3cm}
\textbf{Metas especificas:}
\begin{itemize}
  \item Relacionar datos historicos de ventas con un modelo diferencial.
  \item Resolver numericamente usando Euler, RK2 y RK4.
  \item Comparar precision (error) y eficiencia (tiempo de calculo).
  \item Interpretar resultados para apoyar decisiones empresariales.
\end{itemize}
\end{frame}

\begin{frame}{Reto 1: Datos historicos y modelo}
\textbf{Contexto empresarial:} ventas mensuales (miles de unidades) de un producto en expansion.

\vspace{0.2cm}
Se propone un crecimiento con saturacion de mercado:
\[
\frac{dY}{dt} = rY\left(1 - \frac{Y}{K}\right),
\]
donde:
\begin{itemize}
  \item $Y(t)$: ventas acumuladas o nivel de ventas,
  \item $r$: tasa de crecimiento,
  \item $K$: capacidad maxima del mercado.
\end{itemize}

\textbf{Ajuste de parametros:} se estiman $r$ y $K$ por minimizacion de error cuadratico contra datos historicos.

\vspace{0.2cm}
\textbf{Resultados del ajuste ejecutado en Python:}
\[
r = 0.24, \qquad K = 980, \qquad SSE = 1838.33
\]
\end{frame}

\begin{frame}{Reto 2: Modelo diferencial de primer orden}
\textbf{EDO seleccionada:} modelo logistico (primer orden, no lineal).

\vspace{0.3cm}
\textbf{Solucion exacta (referencia):}
\[
Y(t)=\frac{K}{1+\left(\frac{K-Y_0}{Y_0}\right)e^{-rt}}.
\]

\textbf{Interpretacion empresarial:}
\begin{itemize}
  \item Al inicio, las ventas crecen rapido por adopcion del producto.
  \item Con el tiempo, la demanda se estabiliza por saturacion de clientes.
  \item El valor $K$ permite planificar capacidad, inventarios y metas comerciales.
\end{itemize}
\end{frame}

\begin{frame}{Reto 3: Implementacion numerica en Python}
\textbf{Metodos implementados:}
\begin{itemize}
  \item Euler explicito.
  \item Runge-Kutta de 2\textdegree orden (Heun/trapecio explicito).
  \item Runge-Kutta de 4\textdegree orden.
\end{itemize}

\textbf{Script principal:} \texttt{reto\_metodos\_numericos.py}

\vspace{0.2cm}
\textbf{Salidas generadas:}
\begin{itemize}
  \item \texttt{trayectorias\_ventas.csv} (trayectorias numericas y exacta).
  \item \texttt{benchmark\_metodos.csv} (error final y tiempo medio).
\end{itemize}
\end{frame}

\begin{frame}[fragile]{Fragmento de codigo (Python)}
\begin{lstlisting}
def rk4_step(f, t, y, h):
    k1 = f(t, y)
    k2 = f(t + 0.5*h, y + 0.5*h*k1)
    k3 = f(t + 0.5*h, y + 0.5*h*k2)
    k4 = f(t + h, y + h*k3)
    return y + (h/6.0)*(k1 + 2*k2 + 2*k3 + k4)

# EDO de ventas logistic
# dY/dt = r*Y*(1 - Y/K)
def sales_ode(t, y):
    return r_hat * y * (1.0 - y / k_hat)
\end{lstlisting}
\end{frame}

\begin{frame}{Analisis de error de aproximacion}
\textbf{Metricas usadas:}
\begin{itemize}
  \item MAE: error absoluto medio.
  \item RMSE: raiz del error cuadratico medio.
  \item Error final: $|Y_N - Y(t_f)|$.
\end{itemize}

\vspace{0.2cm}
\textbf{Resultado esperado teorico:}
\begin{itemize}
  \item Euler: orden global $\mathcal{O}(h)$.
  \item RK2: orden global $\mathcal{O}(h^2)$.
  \item RK4: orden global $\mathcal{O}(h^4)$.
\end{itemize}

\textbf{Interpretacion:} para el mismo paso $h$, RK4 suele dar mayor precision; Euler requiere pasos mas pequenos para aproximar igual.
\end{frame}

\begin{frame}{Resultados de error (h = 1.0)}
\begin{table}
\centering
\small
\begin{tabular}{lcccc}
\toprule
Metodo & MAE vs exacta & RMSE vs exacta & MAE vs historico & RMSE vs historico \\
\midrule
Euler & 15.4751 & 17.6279 & 15.9131 & 20.0366 \\
RK2   & 1.1831  & 1.3731  & 10.2896 & 11.8135 \\
RK4   & 0.0020  & 0.0022  & 10.0863 & 11.8912 \\
\bottomrule
\end{tabular}
\end{table}

\vspace{0.1cm}
\textbf{Orden observado (convergencia):}
Euler $\approx 1$, RK2 $\approx 2$, RK4 $\approx 4$.
\end{frame}

\begin{frame}{Analisis de eficiencia computacional}
\textbf{Metodologia:}
\begin{itemize}
  \item Se evaluan varios pasos: $h = 1.0, 0.5, 0.25, 0.125$.
  \item Se mide tiempo promedio por simulacion (ms).
  \item Se compara error final vs solucion exacta.
\end{itemize}

\vspace{0.2cm}
\textbf{Benchmark (resultados medidos):}
\begin{table}
\centering
\begin{tabular}{lccc}
\toprule
Metodo & $h$ & Error final & Tiempo medio (ms) \\
\midrule
Euler & 1.000 & 17.2269 & 0.001856 \\
RK2   & 1.000 & 1.9889  & 0.002706 \\
RK4   & 1.000 & 0.0029  & 0.005720 \\
\bottomrule
\end{tabular}
\end{table}
\end{frame}

\begin{frame}[fragile]{Reproducibilidad (Python)}
Ejecutar en terminal dentro de la carpeta del proyecto:

\begin{lstlisting}
python reto_metodos_numericos.py
\end{lstlisting}

El script imprime tablas de error, orden de convergencia y tiempos, y ademas genera:
\begin{itemize}
  \item \texttt{trayectorias\_ventas.csv}
  \item \texttt{benchmark\_metodos.csv}
\end{itemize}
\end{frame}

\begin{frame}{Discusion empresarial}
\begin{itemize}
  \item El modelo permite anticipar cuando las ventas empiezan a desacelerar.
  \item Facilita decisiones de produccion y logistica al estimar la saturacion del mercado.
  \item Para analisis rapido, Euler puede ser util.
  \item Para planeacion mas fina y proyecciones, RK4 entrega mayor confiabilidad numerica.
\end{itemize}
\end{frame}

\begin{frame}{Conclusiones}
\begin{enumerate}
  \item Se conectaron datos historicos con una EDO de negocio (crecimiento log\'istico).
  \item Se resolvio numericamente con Euler, RK2 y RK4 en Python.
  \item Se evidencio la relacion precision-costo computacional entre metodos.
  \item RK4 fue el mas preciso; Euler el mas simple; RK2 un compromiso intermedio.
\end{enumerate}
\end{frame}

\begin{frame}{Trabajo futuro}
\begin{itemize}
  \item Extender a sistema de EDOs (ventas, inventario y backlog).
  \item Incluir estacionalidad y choques de demanda.
  \item Integrar optimizacion de parametros con datos reales de la empresa.
\end{itemize}
\end{frame}

\begin{frame}
\centering
\Huge Gracias
\end{frame}

\end{document}
